\subsection{Hardware de baixo consumo e Eficiência em IoT}

O desenvolvimento de nós de sensoriamento para cidades inteligentes exige hardware que equilibre poder de processamento e autonomia energética.

Como destacado por \citeonline{Sari2022} e \citeonline{Ahmad2021}, o ESP32 destaca-se nesse cenário como uma plataforma robusta de baixo custo, integrando conectividade sem fio e múltiplos modos de operação. Diferente de arquiteturas mais simples, o ESP32 permite a execução de tarefas complexas e segurança criptográfica sem comprometer severamente o ciclo de vida da bateria, desde que devidamente configurado.

A viabilidade de dispositivos IoT em campo, como medidores de água residenciais, dependem da minimização do consumo em estado de repouso (\textit{deep sleep}). 

Segundo \citeonline{Ferreira2020}, a vida útil da bateria é influenciada criticamente não apenas pelo hardware, mas pela eficiência das rotinas de escrita em periféricos e pela otimização do firmware.

Segundo \citeonline{Camargo2022}, em torno de 60\% do recurso energético de sistemas IOT são consumidos na comunicação, a escolha efetiva de protocolo de comunicação e otimização de parâmetros como dispersão de sinal, largura de banda e taxa de codificação (destinado a redundância de mensagem, tolerância a erros), entre outros fatores, torna-se uma tarefa de suma importância no desenvolvimento de um sistema IOT otimizado de baixo custo.

O \gls{lorawan}, é um protocolo \gls{lpwan}, que possui o objetivo de viabilizar a comunicação em grandes distâncias, conciliado ao baixo consumo energético. Diante disso, o protocolo surge como uma proposta de comunicação de fácil implementação e que propõe uma solução para a problemática de redução de consumo energético, em contra partida, a baixa largura de banda oferecida por tal solução apoia o desenvolvimento de sistemas baseado em edge-computing, trafegando em rede apenas os dados e metadados prontos para serem tratados em gateway central.

Toda via, surge o desafio de alinhar o menor consumo de recurso energético possível, com o objetivo de prolongar a vida útil de bateria dos dispositivos utilizados, às mais otimizadas soluções de rede, hardware e técnicas de desenvolvimento de software.

\subsection{Modelagem OCR}

O \gls{ocr} é o processo de digitalização de dados presentes em imagens e vídeos. A utilização de tal artifício neste trabalho justifica-se pela necessidade de capacitar o sistema em realizar leitura de valores em imagens extraídas pelos nós sensoriais.

Utilizando a metodologia de processamento na borda (TinyML), será adotado o framework TensorFlow Lite for Microcontrollers, escolhido por sua forte integração com a arquitetura do ESP32 e pela otimização elástica no uso de memória \citeonline{david2021}. O pipeline de trabalho consistirá na criação e tratamento do dataset utilizando como base o conjunto público Water Meters Dataset por \citeonline{kucev2020}, associado a capturas locais do próprio sensor, treinamento supervisionado e aplicação do modelo no sistema IoT. O tratamento do dataset incluirá o redimensionamento das imagens, a conversão para escala de cinza, a normalização dos dados e técnicas de data augmentation, visando adequar as entradas às restrições de hardware e melhorar a generalização do aprendizado.

Para o reconhecimento dos dígitos dos hidrômetros, será desenvolvida uma rede neural convolucional (CNN) própria submetida à abordagem de \textit{Transfer Learning} para especialização no cenário estático do projeto. Em complemento, utilizando a abordagem de aprendizagem supervisionada apresentada por \citeonline{SilvaVilela2021}, algoritmos de árvores de decisão e floresta aleatória auxiliarão na classificação do objeto alvo. Tendo a imagem tratada como input, o modelo pré-treinado passará por uma quantização para inteiros de 8 bits (INT8) como apresentado por \citeonline{krishnamoorthi2018}, comprimindo o arquivo final para reduzir o consumo de memória e garantir a execução com fluidez em ambiente bare-metal, de onde espera-se obter como resposta os valores exatos de medição.

Diante da metodologia de processamento em cloud, será aplicado o framework de \gls{ocr} de propósito geral \textit{Tesseract OCR}. Como \citeonline{Gribomont2020} e \citeonline{lederhans2022tratamento} detalham, o Tesseract destaca-se como uma opção de médio custo de processamento e baixa taxa de erro. Além disso, tal framework apresenta alta viabilidade e acessibilidade por ser uma ferramenta open-source, oferecendo vantagens em projetos acadêmicos quando comparada a soluções comerciais pagas, como a Google Vision API.

\subsection{Software Otimizado em Sistemas Restritos}

O desenvolvimento do software para dispositivos IoT exige um rigoroso controle sobre a alocação de recursos, uma vez que a eficiência impacta diretamente ao consumo energético e viabilidade do sistema. Em plataformas com recursos limitados, como o ESP32, as escolhas da linguagem de programação e frameworks são essenciais e determinantes para o sucesso de implementação de tarefas complexas.

Do ponto de vista computacional e eficiência de recursos, o uso nativo do \gls{espidf} (\textit{Espressif IoT Development Framework}) apresenta vantagens significativas em comparação a frameworks de alto nível, como Arduino e MicroPython. O \gls{espidf} é projetado para aproveitar ao máximo as capacidades do hardware, permitindo controle granular sobre a memória e consumo energético, enquanto muitos de seus concorrentes implementam abstrações, podendo ocasionar o uso ineficiente dos limitados recursos do dispositivo \cite{espidf}.

O processamento intensivo exigido pelo \gls{ocr} requer que a implementação evite qualquer desperdício computacional. \citeonline{Plauska2023} demonstram que linguagens compiladas como C e C++ garantem um desempenho superior na gestão de CPU e memória SRAM em relação às linguagens interpretadas. Aliando essa eficiência à total compatibilidade do C/C++ com o ecossistema do \gls{espidf}, o desenvolvedor obtém um controle de memória altamente preciso. Essa gestão minuciosa e de baixo nível é o fator crítico que viabiliza a execução de algoritmos de extração de características em microcontroladores limitados a poucas centenas de kilobytes de memória.

\subsection{Envio de Imagens via LoRaWan}

Uma das principais incógnitas levantadas por este trabalho é a viabilidade 
do envio das imagens geradas em nós na borda da rede ao gateway (nó 
central). Como apresentado por \citeonline{almeida2022iot}, a propagação 
de dados não estruturados via rede \gls{lorawan} é algo viável, mas com a 
ressalva de limitações de bits por pacote, maior latência na 
disponibilização dos dados e a necessidade de algoritmos que orquestrem 
todo o recebimento e envio de dados pela rede, incluindo a fragmentação 
da imagem em múltiplos pacotes e sua posterior remontagem no destino.

Diante disso, casa-se perfeitamente a implementação de algoritmos para 
a redução do tamanho de imagens, a considerar que este é um dos fatores 
mais influentes na execução de tal abordagem. Um ponto apresentado por 
\citeonline{almeida2022iot} é a aplicação de algoritmos de compressão 
como o JPEG e a transformação de imagens em escala de cinza, reduzindo 
assim as informações a serem enviadas via rede.

\subsubsection{Restrições Regulatórias e Parâmetros de Transmissão LoRaWan}

O \textit{Spreading Factor} (\gls{sf}) é um parâmetro central do protocolo 
\gls{lora} que determina a taxa de chirp utilizada na modulação do sinal. 
Valores mais altos de \gls{sf} aumentam a sensibilidade do receptor e o 
alcance da transmissão, porém elevam proporcionalmente o tempo que o 
sinal permanece no ar (\textit{time on air}), impactando diretamente 
o consumo energético e a capacidade da rede~\cite{adam2022}.

O \textit{duty cycle} é a relação entre o tempo em que um dispositivo 
ocupa um canal de rádio e o tempo total de observação. No plano europeu 
EU868 (863--870~MHz), regulações impõem um \textit{duty cycle} máximo 
de 1\%, o que equivale a um orçamento de transmissão de apenas 
36~segundos por hora, ou 864~segundos por dia, conforme 
\citeonline{jebril2018overcoming}.

O \textit{dwell time} é o tempo máximo de ocupação contínua de uma 
frequência por transmissão. No Brasil, o padrão regulamentado pela 
ANATEL para redes \gls{lorawan} é o AU915, operando na faixa de 
902--928~MHz. Diferentemente do EU868, o AU915 não impõe restrição 
de \textit{duty cycle}; em seu lugar, o Ato nº~14.448, de 4 de dezembro 
de 2017~\cite{anatel2017} estabelece que, para sistemas de salto em 
frequência operando nas faixas de 902--907,5~MHz e 915--928~MHz com 
largura de faixa de canal inferior a 250~kHz, caso do \gls{lora} com 
125~kHz, o tempo médio de ocupação de qualquer radiofrequência não 
deve ser superior a \textbf{0,4~segundos} em um intervalo de 
14~segundos, com uso mínimo de 35~canais de salto. O AU915 utiliza 
64~canais de \textit{uplink}, satisfazendo amplamente essa condição.

O tempo no ar por pacote ($T_{\text{pacote}}$) é calculado a partir 
do \gls{sf}, da largura de banda (BW), do coding rate (\gls{cr}) e do 
tamanho do \textit{\gls{payload}}, conforme a formulação oficial do \gls{lora}, 
detalhada nas Equações~\ref{eq:tsym} a \ref{eq:tpacote} da seção de 
Desenvolvimento. O \textit{coding rate} expressa a proporção de bits 
redundantes adicionados para tolerância a erros; o valor 4/5 representa 
o menor \textit{overhead} de redundância disponível no padrão 
\gls{lora}~\cite{adam2022}.

A \textit{Packet Delivery Ratio} (PDR) quantifica a robustez da 
transmissão via \gls{lorawan} pela razão entre o número de pacotes recebidos 
pelo \textit{gateway} e o número de pacotes enviados pelo nó sensor, 
conforme a Equação~\ref{eq:pdr}~\cite{bravo2025ants}.
